\documentclass[aspectratio=169, 10pt]{beamer}
\usepackage{sunass_estilo}

\title{Módulo V: Monopolio y Regulación}
\subtitle{Guía Teórica de Repaso Rápido · PEU SUNASS (Núcleo Regulatorio)}
\author[PEU SUNASS · Economía]{Preparación Examen de Admisión PEU SUNASS}
\institute{Basado en el Temario Oficial de Economía · Microeconomía y Regulación}
\date[2026]{2026}

\begin{document}

\begin{frame}[plain]
  \titlepage
\end{frame}

% ==============================================================================
% SLIDE 2: MAPA CONCEPTUAL DEL MÓDULO V
% ==============================================================================
\begin{frame}{Mapa Conceptual: Monopolio y Regulación}
  \vspace{0.15cm}
  \begin{columns}[T]
    \begin{column}{0.32\textwidth}
      \tarjeta{azulPrincipal}{1. Monopolio Maximizador}{
        Único oferente, barreras a la entrada y demanda decreciente (\(P > IMg\)).\\
        Equilibrio: \(IMg = CMg\), Índice de Lerner y peso muerto (\(PIE\)).
      }
    \end{column}
    \begin{column}{0.32\textwidth}
      \tarjeta{naranjaAcento}{2. Discriminación de Precios}{
        Captura del excedente del consumidor:\\
        1er grado (perfecta), 2do grado (bloques), 3er grado (segmentación) y tarifa en dos partes.
      }
    \end{column}
    \begin{column}{0.32\textwidth}
      \tarjeta{verdeAcento}{3. Monopolio Natural (SUNASS)}{
        Subaditividad de costos (\(CMg < CMe\)).\\
        Regulación tarifaria: \(P=CMe\), Ramsey-Boiteux, Tasa de Retorno vs. Price-Cap (\(RPI-X\)).
      }
    \end{column}
  \end{columns}

  \vspace{0.2cm}
  \recuadroclave{
    \small\textbf{Fórmula Fundamental del Monopolista:} \quad \(IMg(Q) = P(Q)\left(1 - \frac{1}{|\varepsilon_d|}\right)\) \quad \(\implies\) \quad \textbf{Índice de Lerner:} \ \(L = \frac{P - CMg}{P} = \frac{1}{|\varepsilon_d|}\)
  }
  \mensajefinal{Un monopolista maximizador de beneficios con \(CMg > 0\) jamás opera en el tramo inelástico de la demanda (\(|\varepsilon_d| < 1\)).}
\end{frame}

% ==============================================================================
% SLIDE 3: EQUILIBRIO DEL MONOPOLIO, MARK-UP, LERNER Y PÉRDIDA DE EFICIENCIA
% ==============================================================================
\begin{frame}{1. Equilibrio del Monopolista, Mark-Up y Eficiencia}
  \begin{columns}[T]
    \begin{column}{0.52\textwidth}
      {\small\textbf{Equilibrio \((Q_M, P_M)\), Beneficios (\(BT\)) y \(PIE\)}}
      \vspace{0.05cm}
      \begin{tikzpicture}[scale=0.77]
        \begin{axis}[
          axis lines=left,
          xlabel={Cantidad ($Q$)}, ylabel={Precio / Costo},
          xmin=0, xmax=10.5, ymin=0, ymax=12,
          xtick={4, 5.5, 6.67, 11}, xticklabels={$Q_M$, $\frac{a}{2b}$, $Q_C$, $\frac{a}{b}$},
          ytick={3, 4, 7, 11}, yticklabels={$CMg^*$, $CMe^*$, $P_M$, $a$},
          width=8.2cm, height=5.3cm, clip=false
        ]
          % Áreas: CT (azul claro), BT (verde), EC (azul suave), PIE (rojo)
          \addplot[fill=grisTexto!12, draw=none] coordinates {(0,0) (4,0) (4,4) (0,4)} \closedcycle;
          \addplot[fill=verdeAcento!22, draw=none] coordinates {(0,4) (4,4) (4,7) (0,7)} \closedcycle;
          \addplot[fill=azulClaroFondo, draw=none] coordinates {(0,7) (4,7) (0,11)} \closedcycle;
          \addplot[fill=rojoAcento!28, draw=none] coordinates {(4,3) (6.67,4.33) (4,7)} \closedcycle;

          % Demanda P = 11 - Q -> (0,11) a (10,1)
          \addplot[azulPrincipal, very thick, domain=0:10] {11 - x};
          \node[azulPrincipal, font=\scriptsize, anchor=west] at (axis cs:8.3, 2.9) {$D = IMe$};
          % IMg = 11 - 2Q
          \addplot[naranjaAcento, very thick, domain=0:5.5] {11 - 2*x};
          \node[naranjaAcento, font=\scriptsize, anchor=north east] at (axis cs:5.4, 1.4) {$IMg$};
          % CMg = 1 + 0.5Q
          \addplot[verdeAcento, very thick, domain=0:9.8] {1 + 0.5*x};
          \node[verdeAcento, font=\scriptsize, anchor=west] at (axis cs:9.6, 5.9) {$CMg$};
          % CMe = 8/Q + 1 + 0.25Q (mín en Q=5.66)
          \addplot[aceroAcento, thick, dashed, domain=1.15:9.8, samples=50] {8/x + 1 + 0.25*x};
          \node[aceroAcento, font=\scriptsize, anchor=south west] at (axis cs:8.4, 4.15) {$CMe$};

          % Guías y puntos exactos (QM=4, PM=7, CMe=4, CMg=3, QC=6.67)
          \draw[dotted, grisTexto] (axis cs:4, 0) -- (axis cs:4, 7) -- (axis cs:0, 7);
          \draw[dotted, grisTexto] (axis cs:4, 4) -- (axis cs:0, 4);
          \draw[dotted, grisTexto] (axis cs:4, 3) -- (axis cs:0, 3);
          \draw[dotted, grisTexto] (axis cs:6.67, 0) -- (axis cs:6.67, 4.33);
          \addplot[only marks, mark=*, mark size=2pt, azulPrincipal] coordinates {(4,7) (4,4) (4,3) (6.67,4.33)};

          % Etiquetas de áreas
          \node[azulPrincipal, font=\tiny\bfseries] at (axis cs:0.85, 8.85) {$EC$};
          \node[verdeAcento!90!black, font=\tiny\bfseries] at (axis cs:2.05, 6.15) {$BT=\pi_M$};
          \node[grisTexto, font=\tiny\bfseries] at (axis cs:1.8, 1.35) {$CT$};
          \node[rojoAcento, font=\scriptsize\bfseries] at (axis cs:4.85, 4.75) {PIE};
        \end{axis}
      \end{tikzpicture}
    \end{column}

    \begin{column}{0.45\textwidth}
      {\scriptsize\color{grisTexto}Demanda Lineal e Ingreso Marginal}\\[1pt]
      {\small\color{azulPrincipal}\textbf{\(P = a - bQ \implies IMg = a - 2bQ\)}}\\[1pt]
      {\scriptsize
        • Mismo intercepto vertical \(a\) y \textbf{doble pendiente} (\(-2b\)).\\
        • Beneficio total: \(BT = (P_M - CMe^*)\,Q_M\).\\[3pt]
      }
      {\scriptsize\color{grisTexto}Regla Óptima de Mark-Up e Índice de Lerner}\\[1pt]
      {\small\color{naranjaAcento}\textbf{\(P = \frac{CMg}{1 - 1/|\varepsilon_d|} = \frac{CMg}{1 + 1/\varepsilon_d}\)}}\\[1pt]
      {\scriptsize
        • De \(IMg = P(1 + 1/\varepsilon_d) = CMg\) (\(\varepsilon_d < 0\)).\\
        • \textbf{Lerner:} \(L = \frac{P-CMg}{P} = \frac{1}{|\varepsilon_d|} = \frac{-1}{\varepsilon_d} \in [0, 1]\).\\[3pt]
      }
      {\scriptsize\color{grisTexto}Sin Curva de Oferta y Efectos de Bienestar}\\[1pt]
      {\scriptsize
        • Sin relación única \(P \leftrightarrow Q\) (depende de \(\varepsilon_d\)).\\
        • \textbf{Rent-seeking:} gasto improductivo por capturar rentas.\\
        • \textbf{Impuesto \(t\):} con \(\varepsilon_d\) cte., \(\frac{dP}{dt} = \frac{|\varepsilon|}{|\varepsilon|-1} > 1\) (\(>100\%\)).
      }
    \end{column}
  \end{columns}

  \mensajefinal{El monopolio restringe \(Q_M < Q_C\), eleva \(P_M > CMg\) y genera una pérdida irrecuperable (\(PIE\)).}
\end{frame}

% ==============================================================================
% SLIDE 4: PANORAMA DE DISCRIMINACIÓN DE PRECIOS Y TARIFA EN DOS PARTES
% ==============================================================================
\begin{frame}{2. Discriminación de Precios: Panorama Comparativo}
  \vspace{0.08cm}
  \begin{columns}[T]
    \begin{column}{0.32\textwidth}
      \tarjeta{verdeAcento}{1er Grado (Perfecta)}{
        Cobra a cada consumidor/unidad su máxima disposición a pagar.\\[4pt]
        • Extrae \textbf{100\% del \(EC\)} (\(\pi = W\)).\\
        • Produce el nivel eficiente \(Q_C\) \(\implies \mathbf{PIE = 0}\).
      }
    \end{column}
    \begin{column}{0.32\textwidth}
      \tarjeta{naranjaAcento}{2do Grado (Autoselección)}{
        Menús, tarifas por bloques o dos tramos \(T = A + P^*Q\):\\[4pt]
        • El cliente revela su tipo al elegir.\\
        • \textit{Usado en EPS (bloques crecientes IBT y cargo fijo + variable).}
      }
    \end{column}
    \begin{column}{0.32\textwidth}
      \tarjeta{azulPrincipal}{3er Grado (Segmentación)}{
        Precios distintos a grupos observables (doméstico vs. comercial):\\
        \(IMg_1(q_1) = IMg_2(q_2) = CMg\)\\[3pt]
        \(\frac{P_1}{P_2} = \frac{1 - 1/|\varepsilon_2|}{1 - 1/|\varepsilon_1|}\).\\
        \textbf{Mayor precio al grupo más inelástico.}
      }
    \end{column}
  \end{columns}

  \vspace{0.12cm}
  \recuadroclave{
    \small\textbf{Requisitos Esenciales para Discriminar:} \quad \textbf{1)} Poder de mercado (\(P > CMg\)) \quad•\quad \textbf{2)} Conocer o inducir disposición a pagar \quad•\quad \textbf{3)} \textbf{Impedir el arbitraje} (reventa entre clientes).
  }
  \mensajefinal{La discriminación transforma excedente del consumidor en beneficio; la de 3er grado solo eleva el bienestar si expande \(Q\) total.}
\end{frame}

% ==============================================================================
% SLIDE 5: DISCRIMINACIÓN DE PRIMER GRADO (PERFECTA)
% ==============================================================================
\begin{frame}{2a. Discriminación de 1er Grado (Perfecta)}
  \begin{columns}[T]
    \begin{column}{0.52\textwidth}
      {\small\textbf{Extracción Completa del Excedente (\(EC = 0,\; PIE = 0\))}}
      \vspace{0.05cm}
      \begin{tikzpicture}[scale=0.77]
        \begin{axis}[
          axis lines=left,
          xlabel={Cantidad ($Q$)}, ylabel={Precio / Costo},
          xmin=0, xmax=10.5, ymin=0, ymax=11.5,
          xtick={3.6, 6, 10}, xticklabels={$Q_M$, $Q^*=Q_C$, $\frac{a}{b}$},
          ytick={4, 10}, yticklabels={$P^*=CMg^*$, $P_{\max}=a$},
          width=8.2cm, height=5.2cm, clip=false
        ]
          % Área de beneficio bruto total capturado bajo IMe y sobre CMg
          \addplot[fill=verdeAcento!24, draw=none] coordinates {(0,1) (6,4) (0,10)} \closedcycle;

          % Curva de demanda (actúa como IMg efectivo)
          \addplot[azulPrincipal, very thick, domain=0:9.6] {10 - x};
          \node[azulPrincipal, font=\scriptsize, anchor=west] at (axis cs:6.8, 3.5) {$D = IMe = IMg_{1^\circ}$};
          % IMg de precio único (referencia)
          \addplot[naranjaAcento, thick, dashed, domain=0:5.0] {10 - 2*x};
          \node[naranjaAcento, font=\scriptsize, anchor=south west] at (axis cs:3.95, 2.0) {$IMg_{\text{único}}$};
          % CMg = 1 + 0.5Q
          \addplot[verdeAcento, very thick, domain=0:9.6] {1 + 0.5*x};
          \node[verdeAcento, font=\scriptsize, anchor=west] at (axis cs:8.7, 5.5) {$CMg$};

          % Guías y puntos
          \draw[dotted, grisTexto] (axis cs:6, 0) -- (axis cs:6, 4) -- (axis cs:0, 4);
          \draw[dotted, grisTexto!75] (axis cs:3.6, 0) -- (axis cs:3.6, 6.4);
          \addplot[only marks, mark=*, mark size=2.2pt, rojoAcento] coordinates {(6,4)};
          \node[verdeAcento!85!black, font=\tiny\bfseries, align=center] at (axis cs:1.8, 5.5) {Beneficio extraído\\(\(\Delta BT\) sin \(PIE\))};
        \end{axis}
      \end{tikzpicture}
    \end{column}

    \begin{column}{0.45\textwidth}
      {\scriptsize\color{grisTexto}Mecanismo de Precios Individualizados}\\[1pt]
      {\small\color{azulPrincipal}\textbf{\(P^* \le P(q) \le P_{\max} = a\)}}\\[2pt]
      {\scriptsize
        • Cada unidad adicional \(dq\) se vende al precio de reserva del cliente.\\
        • Como vender una unidad más no exige rebajar las anteriores, \textbf{la curva \(IMe\) actúa como \(IMg\) efectivo}.\\[4pt]
      }
      {\scriptsize\color{grisTexto}Optimización y Cálculo de Beneficios}\\[1pt]
      {\small\color{verdeAcento}\textbf{Condición: \(IMe(Q^*) = CMg(Q^*)\)}}\\[2pt]
      {\normalsize\color{azulOscuro}\textbf{\(BT = \int_0^{Q^*} \big[IMe(Q) - CMg(Q)\big]\,dQ - CF\)}}\\[3pt]
      {\scriptsize
        • \textbf{Eficiencia:} produce \(Q^* = Q_C \implies \mathbf{PIE = 0}\).\\
        • \textbf{Equidad:} el productor captura 100\% del bienestar (\(EC = 0\)).\\
        • \textbf{Límite práctico:} exige conocer cada disposición a pagar (asimetría de información).
      }
    \end{column}
  \end{columns}

  \mensajefinal{En 1er grado hay eficiencia plena (\(Q^* = Q_C,\; PIE = 0\)), pero el productor captura todo el excedente.}
\end{frame}

% ==============================================================================
% SLIDE 6: DISCRIMINACIÓN DE SEGUNDO GRADO Y TARIFA EN DOS TRAMOS
% ==============================================================================
\begin{frame}{2b. Discriminación de 2do Grado: Tarifa en Dos Tramos}
  \begin{columns}[T]
    \begin{column}{0.52\textwidth}
      {\small\textbf{Tarifa \(T(Q) = A + P^*Q\) con Dos Demandas Heterogéneas}}
      \vspace{0.05cm}
      \begin{tikzpicture}[scale=0.77]
        \begin{axis}[
          axis lines=left,
          xlabel={Cantidad ($Q$)}, ylabel={Precio / Costo},
          xmin=0, xmax=10.5, ymin=0, ymax=10.8,
          xtick={3, 7}, xticklabels={$Q_2^*$, $Q_1^*$},
          ytick={1, 3, 6, 10}, yticklabels={$CMg$, $P^*$, $B$, $A_{\max}$},
          width=8.2cm, height=5.2cm, clip=false
        ]
          % Cuota de entrada A = EC de la demanda menos dispuesta (IMe_2)
          \addplot[fill=rojoAcento!20, draw=none] coordinates {(0,3) (0,6) (3,3)} \closedcycle;
          % Beneficio variable BV2 & BV1 sobre Q2 y BV1 sobre (Q1 - Q2)
          \addplot[fill=azulPrincipal!20, draw=none] coordinates {(0,1) (3,1) (3,3) (0,3)} \closedcycle;
          \addplot[fill=verdeAcento!22, draw=none] coordinates {(3,1) (7,1) (7,3) (3,3)} \closedcycle;

          % Demanda Alta IMe1: P = 10 - Q
          \addplot[azulPrincipal, very thick, domain=0:9.6] {10 - x};
          \node[azulPrincipal, font=\scriptsize, anchor=west] at (axis cs:6.8, 3.5) {$IMe_1$ (Alta)};
          % Demanda Baja IMe2: P = 6 - Q
          \addplot[rojoAcento, very thick, domain=0:5.8] {6 - x};
          \node[rojoAcento, font=\scriptsize, anchor=south west] at (axis cs:2.1, 4.05) {$IMe_2$ (Baja)};
          % CMg = Cvme = 1 cte.
          \addplot[verdeAcento, very thick, domain=0:9.8] {1};
          \node[verdeAcento, font=\scriptsize, anchor=south west] at (axis cs:8.2, 1.05) {$CMg = C$};

          % Guías y puntos óptimos (P*=3, Q2=3, Q1=7)
          \draw[dotted, grisTexto] (axis cs:3, 0) -- (axis cs:3, 3) -- (axis cs:0, 3);
          \draw[dotted, grisTexto] (axis cs:7, 0) -- (axis cs:7, 3) -- (axis cs:3, 3);
          \addplot[only marks, mark=*, mark size=2pt, azulOscuro] coordinates {(3,3) (7,3)};

          % Etiquetas de regiones
          \node[rojoAcento, font=\scriptsize\bfseries] at (axis cs:0.95, 4.0) {$A$};
          \node[azulPrincipal, font=\tiny\bfseries] at (axis cs:1.5, 2.0) {$BV_1 \land BV_2$};
          \node[verdeAcento!85!black, font=\tiny\bfseries] at (axis cs:5.0, 2.0) {$BV_1$};
        \end{axis}
      \end{tikzpicture}
    \end{column}

    \begin{column}{0.45\textwidth}
      {\scriptsize\color{grisTexto}Cuota de Entrada (\(A\)) y Tarifa de Uso (\(P^*\))}\\[1pt]
      {\small\color{azulPrincipal}\textbf{\(T(Q) = A + P^*Q\)}}\\[2pt]
      {\scriptsize
        • \textbf{Cuota de entrada (\(A\)):} para no excluir al grupo de menor demanda (\(IMe_2\)), el cargo fijo máximo es su excedente:\\
        \quad \(\displaystyle A = \frac{(B - P^*)\,Q_2(P^*)}{2} = EC_2(P^*)\)\\[3pt]
        • \textbf{Beneficio conjunto (\(n = 2\) grupos):}\\
        \quad \(BT = n\,A + (P - CMg)\,(Q_1 + Q_2) - CF\)\\[3pt]
        • \textbf{Optimización (\(\partial BT / \partial P = 0\)):} con grupos heterogéneos conviene fijar \(\mathbf{P^* > CMg}\) (sacrifica parte de \(A\), pero extrae más a \(IMe_1\)).\\[3pt]
        • \textbf{Caso idéntico (\(1\) tipo):} \(P^* = CMg\) y \(A^* = EC(CMg)\).
      }
    \end{column}
  \end{columns}

  \mensajefinal{En 2do grado el consumidor se autoselecciona con tarifas en dos tramos o bloques (IBT en EPS).}
\end{frame}

% ==============================================================================
% SLIDE 7: DISCRIMINACIÓN DE TERCER GRADO Y CASO MULTIPLANTA POR TRAMOS
% ==============================================================================
\begin{frame}{2c. Discriminación de 3er Grado y Multiplanta por Tramos}
  \begin{columns}[T]
    \begin{column}{0.52\textwidth}
      {\small\textbf{Agregación Horizontal \(IMg^T(Q) = CMg^T(Q) = \lambda^*\)}}
      \vspace{0.05cm}
      \begin{tikzpicture}[scale=0.77]
        \begin{axis}[
          axis lines=left,
          xlabel={Cantidad Total ($Q$)}, ylabel={$IMg,\; CMg$},
          xmin=0, xmax=8.3, ymin=0, ymax=9.8,
          xtick={1, 1.5, 3.5, 4.5, 7.5}, xticklabels={$K$, $Y$, $H$, $Q^T$, $L$},
          ytick={3.0}, yticklabels={$\lambda^*$},
          width=8.2cm, height=5.2cm, clip=false
        ]
          % IMg^T por tramos: I (0,9)-(1,7), II (1,7)-(3.5,4), III (3.5,4)-(7.5,0)
          \addplot[verdeAcento, very thick] coordinates {(0,9) (1,7)};
          \node[verdeAcento, font=\tiny\bfseries, anchor=south west] at (axis cs:0.45, 7.9) {Tramo I};
          \addplot[rojoAcento, very thick] coordinates {(1,7) (3.5,4)};
          \node[rojoAcento, font=\tiny\bfseries, anchor=south west] at (axis cs:2.1, 5.4) {Tramo II};
          \addplot[azulPrincipal, very thick] coordinates {(3.5,4) (7.5,0)};
          \node[azulPrincipal, font=\tiny\bfseries, anchor=south west] at (axis cs:5.25, 2.05) {Tramo III (\(IMg^T\))};

          % CMg^T por tramos: I (0,1)-(1.5,2), II (1.5,2)-(3.5,2.8), III (3.5,2.8)-(7.5,3.6)
          \addplot[aceroAcento!75, very thick] coordinates {(0,1) (1.5,2)};
          \addplot[naranjaAcento, very thick] coordinates {(1.5,2) (3.5,2.8)};
          \addplot[rojoAcento!85!black, very thick] coordinates {(3.5,2.8) (7.5,3.6)};
          \node[rojoAcento!85!black, font=\tiny\bfseries, anchor=south] at (axis cs:6.7, 3.55) {\(CMg^T\)};

          % Guías de quiebres y óptimo (Q^T=4.5, lambda*=3.0)
          \draw[dashed, grisTexto!70] (axis cs:1, 0) -- (axis cs:1, 7);
          \draw[dashed, grisTexto!70] (axis cs:1.5, 0) -- (axis cs:1.5, 2);
          \draw[dashed, grisTexto!70] (axis cs:3.5, 0) -- (axis cs:3.5, 4);
          \draw[dotted, azulOscuro, thick] (axis cs:4.5, 0) -- (axis cs:4.5, 3.0) -- (axis cs:0, 3.0);
          \addplot[only marks, mark=*, mark size=2.3pt, azulOscuro] coordinates {(4.5, 3.0)};
          \node[azulOscuro, font=\tiny\bfseries, anchor=south] at (axis cs:4.55, 3.35) {\(IMg^*=CMg^*\)};
        \end{axis}
      \end{tikzpicture}
    \end{column}

    \begin{column}{0.45\textwidth}
      {\scriptsize\color{grisTexto}Algoritmo de Resolución (Segmentos y Plantas)}\\[2pt]
      {\scriptsize
        \textbf{1. Agregar \(IMg^T\):} invertir cada \(IMg_i(q_i)\) y sumar cantidades \(q_i(IMg)\) por tramos según precios de reserva (\(A > B > C\)).\\[2pt]
        \textbf{2. Agregar \(CMg^T\):} invertir cada \(CMg_j(q_j)\) y sumar cantidades \(q_j(CMg)\) por tramos desde la planta más barata.\\[2pt]
        \textbf{3. Óptimo conjunto:} igualar en el tramo activo:\\
        \quad \(\displaystyle IMg^T(Q^T) = CMg^T(Q^T) = \lambda^*\)\\[2pt]
        \textbf{4. Despejar precios y cuotas:} igualar \(\lambda^*\) en cada segmento (\(IMg_i = \lambda^* \Rightarrow q_i^*, P_i^*\)) y planta (\(CMg_j = \lambda^* \Rightarrow q_j^*\)), verificando \(\sum q_i^* = \sum q_j^* = Q^T\).\\[3pt]
      }
      \recuadroclave{
        {\scriptsize\textbf{Regla inversa de elasticidad:} \(\frac{P_1}{P_2} = \frac{1 - 1/|\varepsilon_2|}{1 - 1/|\varepsilon_1|}\)}
      }
    \end{column}
  \end{columns}

  \mensajefinal{En asignación óptima entre segmentos y plantas, todo ingreso y costo marginal activo iguala a \(\lambda^*\).}
\end{frame}

% ==============================================================================
% SLIDE 5: MONOPOLIO NATURAL Y EL DILEMA REGULATORIO
% ==============================================================================
\begin{frame}{3. El Monopolio Natural y el Dilema de la Regulación}
  \begin{columns}[T]
    \begin{column}{0.53\textwidth}
      {\small\textbf{Tarificación First-Best (\(P=CMg\)) vs. Second-Best (\(P=CMe\))}}
      \vspace{0.05cm}
      \begin{tikzpicture}[scale=0.78]
        \begin{axis}[
          axis lines=left,
          xlabel={Cantidad ($Q$)}, ylabel={Tarifa / Costo},
          xmin=0, xmax=10, ymin=0, ymax=10,
          xtick={4.65, 6.8}, xticklabels={$Q_{CMe}$, $Q_{CMg}$},
          ytick={2.2, 4.35}, yticklabels={$P_{CMg}$, $P_{CMe}$},
          width=8.2cm, height=5.2cm
        ]
          % Déficit exacto si P = CMg: CMe(6.8) = 2.2 + 10/6.8 = 3.67
          \addplot[fill=rojoAcento!25, draw=none] coordinates {(0,2.2) (6.8,2.2) (6.8,3.67) (0,3.67)} \closedcycle;
          % CMe exacto con CF=10 y CMg=2.2 cte.: CMe = 2.2 + 10/x
          \addplot[azulPrincipal, very thick, domain=1.3:9.2, samples=50] {2.2 + 10/x};
          \node[azulPrincipal, font=\scriptsize, anchor=south] at (axis cs:8.3, 3.55) {$CMe$};
          % CMg constante = 2.2
          \addplot[verdeAcento, very thick, domain=0:9.5] {2.2};
          \node[verdeAcento, font=\scriptsize, anchor=north] at (axis cs:8.8, 2.1) {$CMg$};
          % Demanda P = 9 - Q -> corta CMe en (4.65, 4.35) y CMg en (6.8, 2.2)
          \addplot[naranjaAcento, very thick, domain=0:9] {9 - x};
          \node[naranjaAcento, font=\scriptsize, anchor=west] at (axis cs:7.6, 1.5) {$D$};

          % Puntos exactos
          \draw[dotted, grisTexto] (axis cs:4.65, 0) -- (axis cs:4.65, 4.35) -- (axis cs:0, 4.35);
          \draw[dotted, grisTexto] (axis cs:6.8, 0) -- (axis cs:6.8, 3.67) -- (axis cs:0, 3.67);
          \addplot[only marks, mark=*, mark size=2.2pt, azulPrincipal] coordinates {(4.65, 4.35)};
          \addplot[only marks, mark=*, mark size=2.2pt, verdeAcento] coordinates {(6.8, 2.2)};
          \node[rojoAcento, font=\tiny] at (axis cs:3.2, 2.95) {\textbf{Déficit si \(P=CMg\)}};
        \end{axis}
      \end{tikzpicture}
    \end{column}

    \begin{column}{0.44\textwidth}
      {\scriptsize\color{grisTexto}Definición de Monopolio Natural}\\[1pt]
      {\large\color{azulPrincipal}\textbf{\(CMe\) decreciente \(\implies CMg < CMe\)}}\\[4pt]
      {\scriptsize
        • Altos costos fijos hundidos (redes de agua/desagüe) y bajos \(CMg\): duplicar redes es socialmente ineficiente (subaditividad).\\[4pt]
      }
      {\scriptsize\color{grisTexto}El Dilema Regulatorio}\\[1pt]
      {\large\color{rojoAcento}\textbf{\(P = CMg \implies \pi < 0\) (Quiebra)}}\\[4pt]
      {\scriptsize
        • \textbf{First-Best (\(P = CMg\)):} Máxima eficiencia (\(PIE=0\)), pero genera pérdida igual al \(CF\) (requiere subsidio estatal).\\
        • \textbf{Second-Best (\(P = CMe\)):} Autofinanciamiento (\(\pi = 0\)), aunque con una pequeña \(PIE\).\\
        • \textbf{Tarifa en Dos Partes:} Cargo fijo = \(CF/N\), Cargo variable = \(CMg\).
      }
    \end{column}
  \end{columns}

  \mensajefinal{En saneamiento (SUNASS), las tarifas suelen diseñarse para cubrir costos medios de largo plazo (autofinanciamiento) más subsidios cruzados.}
\end{frame}

% ==============================================================================
% SLIDE 6: MODELOS DE REGULACIÓN TARIFARIA (RAMSEY, ROR, PRICE-CAP)
% ==============================================================================
\begin{frame}{4. Mecanismos de Regulación Económica (SUNASS)}
  \vspace{0.1cm}
  \begin{columns}[T]
    \begin{column}{0.32\textwidth}
      \tarjeta{azulPrincipal}{Precios de Ramsey-Boiteux}{
        En monopolios multiproducto (agua + alcantarillado), minimiza la \(PIE\) sujeto a \(\pi=0\):\\[3pt]
        \(\frac{P_i - CMg_i}{P_i} = \frac{\lambda}{|\varepsilon_i|}\)\\[3pt]
        \textbf{Regla:} Cargar mayor margen al bien con demanda más \textbf{inelástica}.
      }
    \end{column}
    \begin{column}{0.32\textwidth}
      \tarjeta{rojoAcento}{Tasa de Retorno (Rate of Return)}{
        Fija \(P\) para garantizar un retorno \(s > r\) sobre el capital invertido.\\[3pt]
        • \textbf{Efecto Averch-Johnson:} Incentivo a \textbf{sobrecapitalizar} (\(K/L\) ineficientemente alto).\\
        • Bajo incentivo a reducir costos.
      }
    \end{column}
    \begin{column}{0.32\textwidth}
      \tarjeta{verdeAcento}{Price-Cap (\(RPI - X\))}{
        Fija un techo tarifario ajustado por inflación (\(RPI\)) menos factor de productividad (\(X\)).\\[3pt]
        • \textbf{Alto poder de incentivo:} la firma se queda con el ahorro de costos entre revisiones (5 años en SUNASS).\\
        • Riesgo: degradación de la calidad.
      }
    \end{column}
  \end{columns}

  \vspace{-0.04cm}
  \recuadroclave{
    \small\textbf{Competencia por Comparación (Yardstick Competition):} El regulador usa los costos de otras EPS comparables (``empresa modelo eficiente'') para fijar la tarifa, reduciendo la asimetría de información.
  }
  \mensajefinal{Price-Cap (\(RPI-X\)) y empresa modelo impulsan eficiencia; la tasa de retorno protege el capital pero sobrecapitaliza.}
\end{frame}

% ==============================================================================
% SLIDE 10: MONOPOLIO MULTIPLANTA, COASE Y DOBLE MARGINALIZACIÓN
% ==============================================================================
\begin{frame}{5. Monopolio Multiplanta, Bien Duradero y Doble Marginalización}
  \vspace{0.01cm}
  \begin{columns}[T]
    \begin{column}{0.32\textwidth}
      \tarjeta{azulPrincipal}{Monopolio Multiplanta}{
        Reparte \(Q = q_1 + q_2\) entre plantas igualando costos marginales:\\[2pt]
        \(IMg(Q) = CMg_1(q_1) = CMg_2(q_2)\)\\[2pt]
        • Si \(CMg_1 < CMg_2\) siempre, cierra la planta 2 o la activa cuando \(CMg_1\) alcanza \(\min CMg_2\).
      }
    \end{column}
    \begin{column}{0.32\textwidth}
      \tarjeta{rojoAcento}{Bien Duradero (Coase)}{
        Quien vende bienes duraderos \textbf{compite consigo mismo en el futuro}.\\[2pt]
        • Los clientes esperan rebajas futuras, forzando \(P \to CMg\).\\
        • \textbf{Soluciones:} alquilar (\textit{leasing}), obsolescencia o recompra.
      }
    \end{column}
    \begin{column}{0.32\textwidth}
      \tarjeta{verdeAcento}{Doble Marginalización}{
        Monopolio \textit{upstream} vende a monopolio \textit{downstream}:\\
        • Doble margen: \(P > P_M\), \(Q < Q_M\), menor \(\pi\) total.\\
        • La \textbf{integración vertical} elimina el doble margen y \textbf{baja \(P\)}.
      }
    \end{column}
  \end{columns}

  \vspace{0.02cm}
  \recuadroclave{
    \small\textbf{Impuesto Ad-Valorem vs. Específico en Monopolio:} Para una misma recaudación fiscal, un impuesto \textbf{ad-valorem} (\(\tau\)) induce mayor producción \(Q\) y menor \(P\) que un impuesto específico (\(t\)), por lo que domina en bienestar.
  }
  \mensajefinal{La integración vertical entre monopolios sucesivos beneficia tanto a las firmas como a los consumidores.}
\end{frame}

% ==============================================================================
% SLIDE 11: PEAK-LOAD, SUBSIDIOS CRUZADOS Y CAPTURA REGULATORIA
% ==============================================================================
\begin{frame}{6. Peak-Load Pricing, Subsidios Cruzados y Captura}
  \vspace{0.01cm}
  \begin{columns}[T]
    \begin{column}{0.32\textwidth}
      \tarjeta{azulPrincipal}{Tarifas Peak-Load (Hora Punta)}{
        Demanda fluctúa entre valle (\(D_1\)) y punta (\(D_2\)) con capacidad \(\bar{K}\):\\
        • \textbf{Valle:} \(P_1 = b\) (solo \(CMg\) operativo).\\
        • \textbf{Punta:} \(P_2 = b + \beta\) (\(CMg\) operativo + costo marginal de capacidad \(\beta\)).
      }
    \end{column}
    \begin{column}{0.32\textwidth}
      \tarjeta{naranjaAcento}{Subsidios Cruzados (SUNASS)}{
        Usuarios no residenciales pagan \(P > CMe\) para subsidiar a vulnerables (\(P < CMe\)).\\[2pt]
        • Financia acceso sin aporte fiscal.\\
        • \textbf{Riesgo:} \textit{cream-skimming} (entrada selectiva al segmento rentable) si no hay exclusividad.
      }
    \end{column}
    \begin{column}{0.32\textwidth}
      \tarjeta{rojoAcento}{Captura Regulatoria (Stigler)}{
        Teoría económica de la regulación:\\
        • Grupos concentrados (empresas) tienen mayor incentivo a influir en el regulador que usuarios dispersos.\\
        • Audiencias públicas y transparencia mitigan la captura.
      }
    \end{column}
  \end{columns}

  \vspace{0.04cm}
  \recuadroclave{
    \small\textbf{Empaquetamiento (Bundling):} Vender dos bienes juntos puede elevar beneficios si las valoraciones de los consumidores están \textbf{negativamente correlacionadas}.
  }
  \mensajefinal{Las tarifas diferenciadas por bloque y horario (\textit{peak-load}) asignan eficientemente la capacidad de red.}
\end{frame}

% ==============================================================================
% SLIDE 12: CONCLUSIONES Y CHEAT-SHEET MÓDULO V
% ==============================================================================
\begin{frame}{Conclusiones y Formulario Rápido — Módulo V}
  \vspace{0.04cm}
  \recuadroclave{
    \small\textbf{Regla de Oro:} El monopolista fija \(IMg = CMg\) con poder de mercado \(L = 1/|\varepsilon_d|\); el regulador busca acercar \(P\) hacia \(CMg\) o \(CMe\) preservando la sostenibilidad financiera.
  }

  \vspace{0.12cm}
  \begin{columns}[T]
    \begin{column}{0.32\textwidth}
      \tarjeta{verdeAcento}{Fórmulas Imprescindibles}{
        • Mark-Up: \(P = \frac{CMg}{1 - 1/|\varepsilon_d|}\)\\
        • 2 tramos: \(T = A + P^*Q\) (\(A = EC_2\))\\
        • 3er grado: \(IMg_1 = IMg_2 = CMg\)\\
        • Ramsey: \(\frac{P_i-CMg_i}{P_i} = \frac{\lambda}{|\varepsilon_i|}\)
      }
    \end{column}
    \begin{column}{0.32\textwidth}
      \tarjeta{azulPrincipal}{Trampas de Examen}{
        • Creer que el monopolista tiene curva de oferta.\\
        • Olvidar el Efecto Averch-Johnson (exceso de \(K\) en Tasa de Retorno).\\
        • Confundir 1er grado (\(PIE=0\)) con precio único.
      }
    \end{column}
    \begin{column}{0.32\textwidth}
      \tarjeta{rojoAcento}{Siguientes Pasos}{
        • Practicar Bloque C del Módulo 05 en el Simulador (casos EPS).\\
        • Dominar el cálculo de \(PIE\) y tarifas en dos partes.
      }
    \end{column}
  \end{columns}

  \vspace{0.25cm}
  \begin{center}
    {\color{azulPrincipal}\Large\textbf{¡Éxito!}} \quad {\color{grisTexto}Módulo V · Monopolio y Regulación · PEU SUNASS 2026}
  \end{center}
\end{frame}

\end{document}
